\section{Future Work}
\label{sec:future_work}

\subsection{Scaling Neural Substrates}

The current substrates --- 5,196 SNN neurons, 5,000 LSM neurons, 131,072 HTM cells --- run in real-time on consumer hardware. Scaling to larger substrates (millions of neurons, deeper reservoir architectures, richer HTM column structures) would test whether the assessment results in Section~\ref{sec:results} are scale-dependent. Does integrated information increase with larger SNN networks? Does recurrence quality improve with deeper LSM reservoirs? The modular substrate design allows these experiments without modifying the consciousness modules.

\subsection{Longitudinal Studies}

The most significant gap in our current results is the absence of longitudinal data. Running the framework continuously over days or weeks --- embedded in a system with persistent memory and real input --- would reveal whether assessment scores stabilize, diverge, or exhibit developmental trajectories analogous to biological consciousness development. Of particular interest is whether cross-theory interaction effects (Section~\ref{sec:ablation}) strengthen or change character over extended operation.

\subsection{Embodiment Integration}

The framework's embedding interfaces (Section~\ref{sec:embedding}) are designed but untested in this paper. Connecting the framework to a conversational agent, a robotic platform (physical or simulated), or an interactive simulation would test whether genuine sensory grounding changes the dynamics of the consciousness cycle. The homeostatic drives are the primary connection point: mapping curiosity to information-seeking behavior, social connection to interaction frequency, and novelty to environmental exploration.

\subsection{Additional Theories}

The modular architecture supports adding consciousness theories without disrupting existing modules. Candidate additions include:

\begin{itemize}
    \item \textbf{Predictive Processing in continuous domains} --- implementing Friston's full variational formulation rather than discrete-state active inference
    \item \textbf{Enactivist modules} --- modeling sensorimotor contingencies and embodied sense-making
    \item \textbf{Perceptual Control Theory} \cite{powers1973} --- adding feedback control hierarchies
    \item \textbf{Quantum coherence models} --- if scalable quantum simulation becomes feasible
\end{itemize}

Each addition would introduce new assessment indicators and potentially new cross-theory interactions.

\subsection{Improved IIT Measurement}

Our current $\Phi$ approximation is limited by PyPhi's incompatibility with Python 3.10+. Implementing a native IIT 4.0 calculator --- or contributing Python 3.11+ compatibility to PyPhi --- would enable exact integrated information measurement for small subsystems and better approximations for the full network. This would strengthen the two IIT indicators that currently score lowest.

\subsection{Community Contributions}

The framework is released under Apache 2.0 specifically to invite contribution. We identify several opportunities:

\begin{itemize}
    \item \textbf{Alternative scoring functions} for existing indicators --- testing whether different operationalizations produce different results
    \item \textbf{New ablation configurations} --- testing all 128 possible module combinations
    \item \textbf{Noise normalization refinement} --- developing better baseline noise models
    \item \textbf{Cross-framework benchmarking} --- running the same indicators on other consciousness implementations for comparison
    \item \textbf{Assessment visualization tools} --- interactive dashboards for exploring indicator scores, ablation effects, and longitudinal trajectories
\end{itemize}

We welcome contributions that challenge our implementation choices as much as those that extend them.
